\documentclass[12pt]{article}
\usepackage{amsmath,amssymb,amsthm,mathtools}
\usepackage{fullpage}
\usepackage{hyperref}

\def\C{\mathbb{C}}
\def\N{\mathbb{N}}
\def\R{\mathbb{R}}
\def\E{\mathbb{E}}
\def\P{\mathbb{P}}
\def\ddc{dd^c}
\DeclareMathOperator{\supp}{supp}

\newtheorem{theorem}{Theorem}[section]
\newtheorem{lemma}{Lemma}[section]
\newtheorem{corollary}{Corollary}[section]
\newtheorem{remark}{Remark}[section]

\title{Convergence of the empirical root measure of random polynomial
systems:\\ completing the argument}
\author{Note for the ProbPoly project}
\date{July 2026}

\begin{document}
\maketitle

\noindent
This note supplies the convergence proof that is missing from the section
\emph{Random systems of polynomials} of the working notes. There, the limit of the
empirical root measure is identified by computing, for each coordinate, the conditional
limit distribution given the remaining coordinates, and then assembling the joint limit.
As the notes themselves observe, that argument shows only that \emph{if} the random
measures converge, the limit is the uniform measure on the torus; the convergence itself
(the interchange of limit and integral marked $(*)$) is never established, and the
conditioning step is delicate because at a root the coordinates are functions of the same
random coefficients. Rather than patching that argument, we replace it by a
potential-theoretic one, in the spirit of Shiffman--Zelditch and
Bloom--Shiffman~\cite{BS07}: the empirical root measure is the wedge product of the
normalized zero currents of the $d$ polynomials, each of which concentrates around an
explicit deterministic current. The conditional computations of the working notes are
recovered as one-variable slices of the limiting potential (Remark~\ref{rem:slices}).
Closely related equidistribution theorems exist in the literature (\cite{BS07} for the
total-degree model; see also Bayraktar~\cite{Bay16} for ensembles adapted to general
Newton polytopes), so the theorem below should not be considered new; the point of this
note is a self-contained, quantitative proof tailored to the two models used in the
working notes, under assumptions matching~\cite{BN}.

\section{Setup and statement}

Fix $d \geq 1$. For $n \in \N$ consider a random system $p_1 = \dots = p_d = 0$ of
polynomials in $z = (z_1,\dots,z_d) \in \C^d$,
\[
  p_l(z) \;=\; \sum_{c \,\in\, \Lambda_n} a^{(l)}_c\, z^c, \qquad l = 1,\dots,d,
\]
where the index set is either
\[
  \Lambda_n = \{0,\dots,n\}^d \quad \text{(\emph{model P}: separate degree $n$ in each
  variable)}, \qquad \text{or}
\]
\[
  \Lambda_n = \{c \in \N^d : |c| \le n\} \quad \text{(\emph{model T}: total degree $n$)}.
\]

\medskip\noindent
\textbf{Assumption (A).} The coefficients $a^{(l)}_c$ (over all $l$ and $c$) are
independent complex random variables with mean $0$, variance
$\E|a^{(l)}_c|^2 = 1$, and laws that are absolutely continuous with densities on $\C$
bounded by a common constant $M$. (Standard complex Gaussians satisfy this with
$M = 1/\pi$; the densities need not be identical.)

\medskip
Write
\[
  K_n(z) \;=\; \E\,|p_l(z)|^2 \;=\; \sum_{c\in\Lambda_n} |z^c|^2,
\]
so that in model P, $K_n(z) = \prod_{i=1}^d s_n(|z_i|^2)$ with
$s_n(t) = 1 + t + \dots + t^n$. Let $\sigma$ denote the uniform probability measure on
the unit circle $S^1 \subset \C$, and $\sigma^{\otimes d}$ the Haar probability measure
on the unit torus $(S^1)^d \subset \C^d$.

By Lemma~\ref{lem:generic} below, almost surely the system has exactly $N_n$ isolated
simple roots in $\C^d$, where $N_n = d!\,n^d$ in model P (the BKK number of the cube
$[0,n]^d$) and $N_n = n^d$ in model T (B\'ezout). Define the empirical probability
measure
\[
  \bar\mu_n \;=\; \frac{1}{N_n} \sum_{r \,:\, p_1(r)=\dots=p_d(r)=0} \delta_r .
\]

\begin{theorem}\label{thm:main}
Under Assumption (A), in both models, $\bar\mu_n \to \sigma^{\otimes d}$ weakly in
probability as $n \to \infty$. Quantitatively, for every $f \in C^\infty_c(\C^d)$ there
is $C(f,d,M)$ with
\[
  \E\,\bigl|\langle \bar\mu_n , f\rangle - \langle \sigma^{\otimes d}, f\rangle\bigr|
  \;\le\; C(f,d,M)\,\frac{\log n}{n}.
\]
\end{theorem}

\section{Conventions and genericity}

We use the operator $\ddc \coloneqq \frac{i}{\pi}\partial\bar\partial$, normalized so
that in one variable $\ddc u = \frac{1}{2\pi}\Delta u \cdot \mathrm{Leb}$ and
$\ddc \log|z| = \delta_0$. By the Poincar\'e--Lelong formula, for a polynomial $p$ the
current $[Z_p] = \ddc \log|p|$ is the current of integration over the zero divisor of
$p$. For locally bounded plurisubharmonic (psh) potentials, wedge products of the
currents $\ddc u$ are defined by Bedford--Taylor theory \cite{BT82}; for the currents
$[Z_{p_l}]$, whose potentials $\log|p_l|$ are unbounded on the divisors, the products
\[
  [Z_{p_1}] \wedge \dots \wedge [Z_{p_q}]
\]
are well defined as soon as the divisors intersect properly (each partial intersection
has the expected codimension); see \cite[Ch.~III, \S3--4]{Dem} for the definition, the
associated Fubini--Stokes integration-by-parts formulas, and King's formula, which for a
proper complete intersection identifies
$[Z_{p_1}]\wedge\dots\wedge[Z_{p_d}] = \sum_r m_r \delta_r$, the sum running over the
isolated common roots with their multiplicities. Throughout, $\omega_0 = \ddc |z|^2$ is
the standard K\"ahler form, and the mass of a positive current $S$ of bidimension
$(k,k)$ on a set $A$ is $\|S\|_A = \int_A S \wedge \omega_0^{k}$ (all our currents are
wedges of closed positive $(1,1)$-currents and smooth positive forms, hence strongly
positive, so this is comparable to the total variation of the coefficients). We use
repeatedly: if $\beta$ is a real smooth $(1,1)$-form with $C^0$-norm $\le M_\beta$ on a
compact $K$, then $-c_d M_\beta\, \omega_0 \le \beta \le c_d M_\beta\, \omega_0$ against
strongly positive currents supported in $K$.

\begin{lemma}\label{lem:generic}
Under Assumption (A), almost surely: the divisors $Z_{p_1},\dots,Z_{p_d}$ intersect
properly, all common roots are simple, none lies on a coordinate hyperplane, and their
number is exactly $N_n$. Consequently, almost surely
$\bigwedge_{l=1}^d [Z_{p_l}] = \sum_r \delta_r$ and $\bar\mu_n$ is a probability
measure.
\end{lemma}

\begin{proof}
Each stated property fails only on a proper algebraic (or constructible) subset of the
coefficient space $\C^{d\,|\Lambda_n|}$: for the count this is Bernstein's theorem
\cite{Ber75} in model P (the BKK bound $d!\,n^d = d!\operatorname{vol}([0,n]^d)$ is
attained unless some facial subsystem has a toric root, a proper algebraic condition;
none of the finitely many exceptional conditions holds generically) and B\'ezout in
model T (all $n^d$ projective intersection points are affine and simple outside a proper
subvariety). Roots on $\{z_i = 0\}$ require $d$ equations in $d-1$ variables to have a
common solution, again a proper algebraic condition. A proper algebraic subset has
Lebesgue measure zero, and by Assumption (A) the law of the coefficient vector is
absolutely continuous.
\end{proof}

\section{Deterministic asymptotics of the potentials}

Define
\[
  \bar u_n(z) \;=\; \tfrac{1}{2n} \log K_n(z), \qquad
  \bar T_n \;=\; \ddc \bar u_n ,
\]
a smooth closed positive $(1,1)$-form ($K_n \ge 1 > 0$ everywhere, and $\bar u_n$ is psh
as a logarithm of a sum of squared moduli of holomorphic monomials). Define the limit
potentials
\[
  V_P(z) = \sum_{i=1}^d \log^+ |z_i| , \qquad
  V_T(z) = \max_{1\le i\le d}\, \log^+ |z_i| = \log^+ \max_i |z_i| .
\]

\begin{lemma}\label{lem:potconv}
In model P (resp.\ T), with $V = V_P$ (resp.\ $V_T$),
\[
  \sup_{z \in \C^d}\, \bigl| \bar u_n(z) - V(z) \bigr| \;\le\; \frac{d \log (n+d)}{2n}
  \;=:\; \delta_n .
\]
\end{lemma}

\begin{proof}
Model P: $\max(1,t^n) \le s_n(t) \le (n+1)\max(1,t^n)$ for $t\ge0$, so
$\bigl|\frac{1}{2n}\log s_n(|z_i|^2) - \log^+|z_i|\bigr| \le \frac{\log(n+1)}{2n}$ for
each $i$; sum over $i$. Model T: every monomial obeys
$|z^c|^2 \le \max\bigl(1,\max_i|z_i|\bigr)^{2n}$, and the terms $c = 0$ and $c = n e_i$
give $K_n \ge \max\bigl(1, \max_i |z_i|^{2n}\bigr)$, while
$|\Lambda_n| = \binom{n+d}{d} \le (n+d)^d$.
\end{proof}

\begin{lemma}\label{lem:limitMA}
$\frac{1}{d!} (\ddc V_P)^{\wedge d} = \sigma^{\otimes d}$ and
$(\ddc V_T)^{\wedge d} = \sigma^{\otimes d}$.
\end{lemma}

\begin{proof}
For $V_P$: $\ddc \log^+|z_i|$ is, in one variable, the measure $\sigma(z_i)$; as a
$(1,1)$-current on $\C^d$ it has only a $dz_i \wedge d\bar z_i$ component. For a sum
$u = \sum_i u_i(z_i)$ of locally bounded psh functions of one variable each, smoothing
each $u_i$ separately and using Bedford--Taylor convergence \cite{BT82},
$(\ddc u)^{\wedge d} = d!\, \bigwedge_i \ddc u_i$, because in the expansion of the
$d$-fold wedge every term repeating an index vanishes and the $d!$ orderings of
$\{1,\dots,d\}$ each contribute $\ddc u_1 \wedge \dots \wedge \ddc u_d$, i.e.\ the
product measure. Hence $(\ddc V_P)^{\wedge d} = d!\, \sigma^{\otimes d}$.

For $V_T$: $V_T$ is the pluricomplex Green function (Siciak--Zaharjuta extremal
function) of the closed unit polydisk, and its Monge--Amp\`ere measure is the
equilibrium measure of the polydisk, namely the Haar probability measure on the
distinguished boundary $(S^1)^d$; this is classical, see \cite[\S5.1]{Kli} (also
\cite{BT82}).
\end{proof}

\section{Pointwise anti-concentration estimate}

For $l = 1,\dots,d$ set
\[
  u_l \;=\; \tfrac1n \log |p_l|, \qquad
  w_l \;\coloneqq\; u_l - \bar u_n \;=\; \tfrac{1}{2n} \log \frac{|p_l|^2}{K_n},
\]
so that $\frac1n [Z_{p_l}] = \bar T_n + \ddc w_l$.

\begin{lemma}\label{lem:pointwise}
There is $C_0 = C_0(d,M)$ such that for every $z \in \C^d$ with $K_n(z) \neq 0$ (i.e.\
all $z$) and every $l$,
\[
  \E\,|w_l(z)| \;\le\; C_0\, \frac{\log (n+2)}{n}.
\]
\end{lemma}

\begin{proof}
Fix $z$ and write $\xi = p_l(z)/\sqrt{K_n(z)} = \sum_{c} b_c\, a^{(l)}_c$ with
$b_c = z^c / \sqrt{K_n(z)}$, $\sum_c |b_c|^2 = 1$. Then
$2n\,|w_l(z)| = |\log |\xi|^2| \le \log^+ |\xi|^2 + \log^- |\xi|^2$.

\emph{Upper tail.} By Jensen's inequality,
$\E \log^+ |\xi|^2 \le \E \log (1 + |\xi|^2) \le \log\bigl(1 + \E|\xi|^2\bigr) = \log 2$.

\emph{Lower tail (anti-concentration).} Choose $c^*$ with $|b_{c^*}| = \max_c |b_c|$, so
$|b_{c^*}|^2 \ge 1/|\Lambda_n| \ge (n+d)^{-d}$. Conditioning on all coefficients except
$a^{(l)}_{c^*}$, the conditional law of $\xi$ has a density on $\C$ bounded by
$M/|b_{c^*}|^2$, hence for every $\varepsilon > 0$,
\[
  \P\bigl( |\xi| \le \varepsilon \bigr) \;\le\; \pi M\, \varepsilon^2\, (n+d)^d .
\]
Therefore
\[
  \E \log^- |\xi|^2
  = \int_0^\infty \P\bigl(|\xi|^2 \le e^{-s}\bigr)\, ds
  \le \int_0^\infty \min\bigl( 1,\; \pi M (n+d)^d e^{-s} \bigr)\, ds
  \le \log\bigl(\pi M (n+d)^d\bigr) + 1 .
\]
Combining, $\E|w_l(z)| \le \frac{1}{2n}\bigl( \log 2 + 1 + \log(\pi M) + d\log(n+d)
\bigr)$.
\end{proof}

Note that only this pointwise-in-$z$ expectation bound is needed; no uniform control of
$w_l$ over $z$, and in particular no control near the zero divisors, enters the
argument.

\section{Mass bounds and the telescoping estimate}

The heart of the proof is the following telescoping device. All currents below are
understood on a fixed open set containing the (compact) support of the test data, and
all wedge products are a.s.\ well defined by Lemma~\ref{lem:generic} and
\cite[Ch.~III]{Dem}.

\begin{lemma}[uniform mass bounds]\label{lem:mass}
Let $\chi \in C^\infty_c(\C^d)$, $\chi \ge 0$. For $0 \le q \le d$ and $L \subseteq
\{1,\dots,d\}$ with $|L| = q$, let $\Phi_n$ be any wedge of $d - q$ factors, each equal
to $\bar T_n$ or to $\omega_0$. Then there is $C_1 = C_1(\chi, d, M)$, independent of
$n$, $L$ and the choice of factors, such that
\[
  \E \int \chi\; \Bigl( \bigwedge_{l \in L} \tfrac1n [Z_{p_l}] \Bigr) \wedge \Phi_n
  \;\le\; C_1 .
\]
\end{lemma}

\begin{proof}
Induction on $q$.

\emph{Base case $q = 0$.} The integrand is a smooth positive $(d,d)$-form. Since the
potentials $\bar u_n$ are psh, smooth, and uniformly bounded on compacts (by
Lemma~\ref{lem:potconv} they converge uniformly to $V$, which is locally bounded), the
Chern--Levine--Nirenberg inequality gives the bound; concretely, integrating by parts
one factor at a time,
$\int \chi\, \ddc \bar u_n \wedge \Phi = \int \bar u_n\, \ddc\chi \wedge \Phi$, and
$\ddc \chi$ is dominated by $c_d \|\chi\|_{C^2}\, \omega_0$ against positive currents,
so each step trades a $\bar T_n$ for an $\omega_0$ at the cost of a factor
$c_d \|\chi\|_{C^2} \sup_{\supp\chi} |\bar u_n|$ (after replacing $\chi$ by a fixed
larger cutoff). After $\le d$ steps only powers of $\omega_0$ remain.

\emph{Induction step.} Let $q \ge 1$, pick $m \in L$, and write
\[
  \bigwedge_{l\in L} \tfrac1n [Z_{p_l}] \wedge \Phi_n
  \;=\; R \wedge \tfrac1n [Z_{p_m}],
  \qquad
  R \coloneqq \bigwedge_{l\in L\setminus\{m\}} \tfrac1n [Z_{p_l}] \wedge \Phi_n .
\]
Since $\tfrac1n [Z_{p_m}] = \bar T_n + \ddc w_m$,
\[
  \int \chi\, R \wedge \tfrac1n [Z_{p_m}]
  \;=\; \int \chi\, R \wedge \bar T_n
  \;+\; \int w_m \, \ddc\chi \wedge R ,
\]
where we integrated by parts (\cite[Ch.~III, \S3]{Dem}; $R$ is closed and positive, $w_m$
is a difference of psh functions locally integrable against the trace of $R$ because
both wedge products $R \wedge \frac1n[Z_{p_m}]$ and $R \wedge \bar T_n$ are well
defined). The first term is handled by the induction hypothesis ($q-1$ random factors,
$\Phi$ replaced by $\Phi \wedge \bar T_n$). For the second, dominate
$\ddc \chi$ by $c_d \|\chi\|_{C^2}\, \omega_0$, so
\[
  \Bigl| \int w_m\, \ddc \chi \wedge R \Bigr|
  \;\le\; c_d \|\chi\|_{C^2} \int_{\supp \chi} |w_m|\; d\bigl( R \wedge \omega_0 \bigr).
\]
Now the key point: $R$ is built from $p_l$ with $l \neq m$ and from deterministic forms,
hence is \emph{independent} of $w_m$. Taking expectations, conditioning on $R$, and
applying Fubini together with Lemma~\ref{lem:pointwise},
\[
  \E \int_{\supp\chi} |w_m| \; d(R\wedge\omega_0)
  \;=\; \E_R \int_{\supp\chi} \E\bigl[|w_m(z)|\bigr]\; d(R\wedge\omega_0)(z)
  \;\le\; C_0 \tfrac{\log(n+2)}{n}\; \E\, \| R \wedge \omega_0 \|_{\supp\chi},
\]
and $\E \|R \wedge \omega_0\|_{\supp\chi}$ is again of the form covered by the induction
hypothesis ($q - 1$ random factors, one more $\omega_0$), after inserting a cutoff
$\tilde\chi = 1$ on $\supp \chi$. Since $\frac{\log(n+2)}{n}$ is bounded, the induction
closes.
\end{proof}

\begin{lemma}[telescoping estimate]\label{lem:telescope}
For every $f \in C^\infty_c(\C^d)$ there is $C_2 = C_2(f,d,M)$ with
\[
  \E\, \Bigl| \Bigl\langle \bigwedge_{l=1}^d \tfrac1n [Z_{p_l}]
  \;-\; \bar T_n^{\wedge d}, \; f \Bigr\rangle \Bigr|
  \;\le\; C_2\, \frac{\log n}{n} .
\]
\end{lemma}

\begin{proof}
Telescope:
\[
  \bigwedge_{l=1}^d \tfrac1n [Z_{p_l}] - \bar T_n^{\wedge d}
  \;=\; \sum_{m=1}^{d} \Bigl( \bigwedge_{l<m} \tfrac1n [Z_{p_l}] \Bigr)
        \wedge \ddc w_m \wedge \bar T_n^{\wedge (d-m)} .
\]
Pair with $f$ and integrate by parts in each summand:
\[
  \Bigl\langle \Bigl( \bigwedge_{l<m} \tfrac1n [Z_{p_l}]\Bigr) \wedge \ddc w_m \wedge
  \bar T_n^{\wedge(d-m)}, f \Bigr\rangle
  \;=\; \int w_m \; \ddc f \wedge \Bigl( \bigwedge_{l<m} \tfrac1n [Z_{p_l}] \Bigr)
  \wedge \bar T_n^{\wedge (d-m)} .
\]
The current on the right involves only $p_l$ with $l < m$, hence is independent of
$w_m$. Dominating $\ddc f$ by $c_d\|f\|_{C^2}\,\omega_0$ and repeating the
Fubini step of Lemma~\ref{lem:mass}, each summand is bounded in expectation by
$c_d \|f\|_{C^2}\; C_0 \frac{\log(n+2)}{n}\; C_1$, with $C_1$ from Lemma~\ref{lem:mass}
applied with a cutoff $\chi = 1$ on $\supp f$.
\end{proof}

\section{Proof of Theorem \ref{thm:main}}

Let $V = V_P$ or $V_T$ according to the model, and let
$\lambda \coloneqq \frac{1}{d!}(\ddc V_P)^{\wedge d} = \sigma^{\otimes d}$, resp.\
$\lambda \coloneqq (\ddc V_T)^{\wedge d} = \sigma^{\otimes d}$
(Lemma~\ref{lem:limitMA}). Fix $f \in C^\infty_c(\C^d)$.

\emph{Deterministic part.} The same telescoping as in Lemma~\ref{lem:telescope}, with
$e_n \coloneqq \bar u_n - V$ in place of $w_m$ and $\ddc V$ in place of
$\frac1n[Z_{p_l}]$, gives
\[
  \bigl| \langle \bar T_n^{\wedge d} - (\ddc V)^{\wedge d}, f \rangle \bigr|
  \;=\; \Bigl| \sum_{m=1}^d \int e_n\, \ddc f \wedge \bar T_n^{\wedge(m-1)} \wedge
  (\ddc V)^{\wedge (d-m)} \Bigr|
  \;\le\; C_3(f,d)\, \delta_n ,
\]
using $\|e_n\|_\infty \le \delta_n$ (Lemma~\ref{lem:potconv}) and uniform local mass
bounds for the mixed wedges $\bar T_n^{\wedge(m-1)} \wedge (\ddc V)^{\wedge(d-m)}$
(Chern--Levine--Nirenberg as in the base case of Lemma~\ref{lem:mass}: all potentials
involved are psh and uniformly bounded on compacts, and $\ddc V$-wedges are defined by
Bedford--Taylor theory since $V$ is locally bounded).

\emph{Random part.} By Lemma~\ref{lem:generic},
$\langle \bar\mu_n, f\rangle = \frac{n^d}{N_n} \bigl\langle \bigwedge_l \frac1n
[Z_{p_l}], f \bigr\rangle$ almost surely, and $\frac{n^d}{N_n} \le 1$. Hence, combining
Lemma~\ref{lem:telescope} with the deterministic part and
$\frac{n^d}{N_n}(\ddc V)^{\wedge d} = \lambda$,
\[
  \E \bigl| \langle \bar\mu_n, f \rangle - \langle \sigma^{\otimes d}, f\rangle \bigr|
  \;\le\; \frac{n^d}{N_n}\Bigl( C_2 \tfrac{\log n}{n} + C_3 \delta_n \Bigr)
  \;\le\; C(f,d,M)\, \frac{\log n}{n} .
\]
This proves the quantitative claim; convergence in probability of
$\langle \bar\mu_n, f\rangle$ for each fixed $f \in C^\infty_c$ follows by Markov's
inequality.

\emph{From vague to weak convergence.} Both $\bar\mu_n$ (a.s.) and $\sigma^{\otimes d}$
are probability measures. Choose $R > 1$ and $g \in C^\infty_c$ with
$0 \le g \le 1$ and $g \equiv 1$ on the closed ball $\bar B_R \supset (S^1)^d$. Then
$\langle \bar\mu_n, g\rangle \to \langle \sigma^{\otimes d}, g \rangle = 1$ in
probability, so $\bar\mu_n(\C^d \setminus B_R) \le 1 - \langle \bar\mu_n, g\rangle \to
0$ in probability: no mass escapes to infinity. Together with convergence on a countable
family of $C^\infty_c$ functions that is dense in $C_c$, this upgrades vague convergence
in probability to weak convergence in probability. \qed

\section{Remarks}

\begin{remark}[relation to the conditional argument in the notes]\label{rem:slices}
The conditional computations in the working notes are exactly the one-variable slices of
$\bar u_n$ and $V$. In model T, for fixed $x_2,\dots,x_d$, the variance profile of the
univariate polynomial in $x_1$ computed there has normalized logarithm converging to
$z \mapsto \max\bigl(\log|z|, \log|x_*|\bigr)$ with $|x_*| = \max(1, |x_2|, \dots,
|x_d|)$ --- which is $V_T(z, x_2, \dots, x_d)$ as a function of $z$. Thus those
computations correctly identify the limit potential; what they could not deliver is the
convergence of the random measures, which is precisely what
Lemmas~\ref{lem:pointwise}--\ref{lem:telescope} supply. The interchange-of-limits step
$(*)$ and the proposed extension of \cite[Lemma B.1]{BN} are no longer needed.
\end{remark}

\begin{remark}[scope of the assumptions]
Gaussian coefficients are nowhere needed; Assumption (A) (independent, unit-variance,
uniformly bounded densities) is used only twice: for genericity
(Lemma~\ref{lem:generic}) and for the anti-concentration bound in
Lemma~\ref{lem:pointwise}. In the Gaussian case $|\xi|^2 \sim \mathrm{Exp}(1)$ in
Lemma~\ref{lem:pointwise} and one gets the slightly better rate $C/n$. The variances
need not all equal $1$ either: the proof goes through verbatim whenever the coefficient
variances $\alpha_c^2$ are such that $\frac{1}{2n}\log \sum_c \alpha_c^2 |z^c|^2$
converges uniformly to some locally bounded $V$; the limit measure is then
$(\ddc V)^{\wedge d}$ suitably normalized. This covers, e.g., Kostlan-type variance
profiles, and makes explicit how the limit depends only on the Newton polytope and the
exponential scale of the variances, in line with \cite{Bay16}.
\end{remark}

\begin{remark}[normalization of the root count]
In model P the number of roots is $d!\,n^d$ (the BKK number of the cube), not $n^d$; for
$d = 2$, two bidegree-$(n,n)$ curves meet in $2n^2$ points. The corollary in the working
notes should normalize the empirical measure accordingly; the limit statement is
unaffected.
\end{remark}

\begin{remark}[almost-sure convergence]
The bound $O(\log n / n)$ is not summable, so Borel--Cantelli gives almost-sure
convergence only along polynomially sparse subsequences (e.g.\ $n_j = j^2$). Full
almost-sure convergence can be obtained by a second-moment refinement of
Lemma~\ref{lem:telescope} or by invoking the variance estimates of \cite{BS07}; since
the application in the working notes only requires convergence in probability, we do not
pursue this.
\end{remark}

\begin{thebibliography}{9}

\bibitem{BN} G.~Barbarino and V.~Noferini,
\emph{The limit empirical spectral distribution of complex matrix polynomials},
Random Matrices: Theory and Applications 11(3), 2250023, 2022.

\bibitem{Bay16} T.~Bayraktar,
\emph{Equidistribution of zeros of random holomorphic sections},
Indiana Univ. Math. J. 65 (2016), 1759--1793.

\bibitem{BT82} E.~Bedford and B.~A.~Taylor,
\emph{A new capacity for plurisubharmonic functions},
Acta Math. 149 (1982), 1--40.

\bibitem{Ber75} D.~N.~Bernstein,
\emph{The number of roots of a system of equations},
Funct. Anal. Appl. 9 (1975), 183--185.

\bibitem{BS07} T.~Bloom and B.~Shiffman,
\emph{Zeros of random polynomials on $\C^m$},
Math. Res. Lett. 14 (2007), 469--479.

\bibitem{Dem} J.-P.~Demailly,
\emph{Complex Analytic and Differential Geometry},
open access book, Institut Fourier, Grenoble.

\bibitem{Kli} M.~Klimek,
\emph{Pluripotential Theory},
London Math. Soc. Monographs 6, Oxford University Press, 1991.

\end{thebibliography}

\end{document}
